\section{实验：定位等待，消除重复工作}
\label{sec:evaluation}
实验遵循同一节奏：先由用户旅程确定问题，再以代码路径提出机制假设，最后选取能区分竞争解释的对照。这里报告项目已留存的 2026 年 10 月实验；其环境、原始记录位置与执行命令集中列于附录~\ref{app:repro}。不同工作负载回答不同问题，不混合计算一个综合“性能提升”。

\subsection{等待策略比轮询周期更重要}
最初服务端通过周期性检查接收工作完成结果。即便业务操作很轻，完成对象仍需等待下一次轮询。Windows 本地、单客户端、两工作线程的 100 次顺序 ping 探测中，旧路径中位延迟为 15.532 ms。改为先入完成队列、再写入原生非阻塞回环唤醒套接字后，中位延迟为 0.126 ms，p95 从 16.163 ms 降至 0.222 ms。这里的改变不是缩短超时，也不是提高系统计时器精度，而是让“存在待处理工作”直接成为就绪事件。

相同问题在终端出现了另一个版本。旧交互循环先等待网络 25 ms，再读取一个 UTF-16 按键，因而粘贴和缓冲输入被错误地序列化为逐字符等待。替换为 FTXUI 事件批处理后，真实 Windows 伪控制台（ConPTY）中的 108 字符 ASCII 突发，中位完成时间从 3410.650 ms 降为 26.817 ms；76 码点中文/emoji 突发从 2410.523 ms 降为 28.062 ms（表~\ref{tab:latency}）。单键 p50 仍约 30.6 ms，说明主要收益是移除与字符数成正比的额外等待，而非改变所有输入的固定显示延迟。

\begin{table}[t]
\caption{不同路径的配对延迟实验。各行只在本行环境内比较；突发值为三次完成时间的中位数，单键值为 78 个观测的 p50。}
\label{tab:latency}
\centering\small
\begin{tabular}{lrr}
\toprule
测量路径 & 原路径 (ms) & 改进路径 (ms)\\
\midrule
服务端完成通知：ping p50 & 15.532 & 0.126\\
服务端完成通知：ping p95 & 16.163 & 0.222\\
终端：ASCII 单键 p50 & 30.844 & 30.642\\
终端：108 字符 ASCII 突发 & 3410.650 & 26.817\\
终端：76 码点中/emoji 突发 & 2410.523 & 28.062\\
\bottomrule
\end{tabular}
\end{table}

终端实验同时要求独立 Bob 协议客户端观察到精确正文，避免把本地回显当作发送成功。旧、新客户端都完成了共同工作负载的九条消息；新路径还通过逐键中文/emoji 和光标编辑，独立对端共验证十三条消息。延迟终点是当前屏幕投影出现预期草稿，不是输出日志中曾出现某个片段。读线程持续排空输出，测试宿主响应 Win32 输入模式协商，避免测量管道误入普通输入路径。

\subsection{扇出优化：共享字节，而非共享可变进度}
每条房间事件在各接收者处拥有相同的编码结果。原实现却为每个接收者重新序列化与压缩，包括重复初始化 LZ 编码表。改进后每事件只编码一次，各输出队列共享不可变字节缓冲区，同时保留各自写偏移。背压仍按每接收者的逻辑字节计费，不因底层共享而放松慢客户端限制。

在 Windows GCC Debug、内存后端、四工作线程、16 客户端各发送 100 条 256 字节正文的匹配 A/B 中，三次试验的吞吐中位数从 651.64 ACK/s 提升至 2652.47 ACK/s，发送至 ACK 的试验 p50 中位数从 24.466 ms 降至 5.568 ms（表~\ref{tab:fanout}）。这一结果对应移除重复计算的机制，也保持原有协议与压缩格式。测量包含 Python 接收端和高可压缩文本，是本地扇出工作负载的结果，不作为 PostgreSQL 容量估计。

\begin{table}[t]
\caption{共享编码的三次匹配扇出试验。每次 1600 条消息；改动仅为编码缓冲区复用。}
\label{tab:fanout}
\centering\small
\begin{tabular}{crrrr}
\toprule
 & \multicolumn{2}{c}{ACK/s} & \multicolumn{2}{c}{发送至 ACK p50 (ms)}\\
试验 & 逐接收者 & 共享 & 逐接收者 & 共享\\
\midrule
1 & 659.50 & 2604.69 & 24.086 & 5.767\\
2 & 647.36 & 2652.47 & 24.639 & 5.492\\
3 & 651.64 & 2694.16 & 24.466 & 5.568\\
\bottomrule
\end{tabular}
\end{table}

\subsection{一个被放弃的解释：普遍启用 TCP\_NODELAY}
对于小消息高延迟，Nagle 与延迟 ACK 的交互是自然假设，但机制存在并不意味着它解释当前观测。早期保留二进制与启用 TCP\_NODELAY 的新二进制比较，看似显示明显差异；然而在同一当前源码上撤回选项后，性能没有回到早期水平。这个反向对照揭示了源码身份混杂。当前源码的启用组与撤回组 p50 分别约 24.5--25.1 ms 和 24.3--24.9 ms，没有带来值得保留的改善。

因此，系统只在完成唤醒通道保留窄范围 TCP\_NODELAY，而不对所有业务连接施加同一策略。后续共享编码的匹配 A/B 则直接检验了重复压缩这一解释。这个过程的价值不是得到一个“最佳 TCP 参数”，而是将待优化对象从传输层猜测重新定位到应用内的工作放大。

\subsection{正确性覆盖：从字节流到真实终端}
\label{sec:validation}
性质测试覆盖分片/合并等价性、UTF-8 与有界解压、JSONL 往返和游标转移；进程测试覆盖认证、租户隔离、双 IPC 客户端的碰撞请求 ID、FIN 后命令排空及断线重连。240 个并发同房间事件由独立观察者检查严格 ID 顺序；20 次发送后立即 FIN 与 32 命令合并后 FIN，检查已提交消息不会因请求者生命周期而消失。

真实终端测试进一步覆盖异步收信时编辑、中文/emoji、窄屏、滚动后显式已读、断线草稿、历史窗口往返与退出恢复。附录~\ref{app:repro} 所标识的跨平台执行在 Windows、Linux、macOS、强制 poll 及 ASan/UBSan 配置完成相应的 13 项测试；PostgreSQL 专用旅程还验证服务重启后的字段一致性，以及连接池会话终止后的原进程恢复。与吞吐数字相比，这些旅程更直接刻画系统是否完成了用户想做的事情。
